%----------------------------------------------------------------------------------------
% NOTATIE EN FIGUREN VAN DE CURSUS STELSELS (2de graad)
% Staat in de preambule en niet in een hoofdstuk: ook de lijsten met oplossingen gebruiken dit.
% Vertrokken van notatie.tex van Functies (grafieken, tabellen) en Eerstegraadsvergelijkingen (de balans).
%----------------------------------------------------------------------------------------
% Bewaarde grafiek, geplaatst door maak-pdf.py: bestand en pagina in figuren-cache/
\newcommand{\grafiekbestand}[2]{\includegraphics[page=#2,trim=3bp 3bp 3bp 3bp]{figuren-cache/#1}}

\providecommand{\R}{\mathbb{R}}
% Absolute waarde: \abs{-3}
\providecommand{\abs}[1]{\lvert #1\rvert}

%----------------------------------------------------------------------------------------
% STELSELS
%----------------------------------------------------------------------------------------
% Een stelsel van twee vergelijkingen met een accolade: \stelsel{2x + y = 8}{x + 3y = 9}
% Met de gelijkheidstekens onder elkaar: \stelsel{2x + y &= 8}{x + 3y &= 9}
\newcommand{\stelsel}[2]{\left\{\renewcommand{\arraystretch}{1.2}\begin{array}{@{}r@{}>{{}}l@{}}#1\\#2\end{array}\right.}
% Drie vergelijkingen (opdrachten met drie onbekenden): \stelseldrie{...}{...}{...}
\newcommand{\stelseldrie}[3]{\left\{\renewcommand{\arraystretch}{1.2}\begin{array}{@{}r@{}>{{}}l@{}}#1\\#2\\#3\end{array}\right.}
% Een determinant van orde 2 (pakket 5): \detm{a}{b}{c}{d}
\newcommand{\detm}[4]{\begin{vmatrix}#1 & #2\\ #3 & #4\end{vmatrix}}

% Zoals stappen, maar met een bredere kolom voor de berekening: voor lange formules
\newenvironment{stappenbreed}{%
  \par\renewcommand{\arraystretch}{1.25}%
  \noindent\begin{tabular}{@{}>{$\displaystyle}p{9.8cm}<{$}@{\hspace{4mm}}>{\RaggedRight\small\itshape\arraybackslash}p{4.9cm}@{}}%
}{\end{tabular}\par}

% Zoals stappen, met eigen kolombreedtes: \begin{stappenvrij}{6.2cm}{3.8cm}. Voor naast een figuur.
\newenvironment{stappenvrij}[2]{%
  \par\renewcommand{\arraystretch}{1.25}%
  \noindent\begin{tabular}{@{}>{$\displaystyle}p{#1}<{$}@{\hspace{3mm}}>{\RaggedRight\small\itshape\arraybackslash}p{#2}@{}}%
}{\end{tabular}\par}

% Tekst of berekening links, figuur rechts: \naast{links}{rechts}
\newsavebox{\naastbox}
\newcommand{\naast}[3][]{%
  \par\noindent\sbox{\naastbox}{#3}%
  \begin{minipage}[t]{\dimexpr\linewidth-\wd\naastbox-3mm\relax}\vspace{0pt}#2\end{minipage}\hfill
  \begin{minipage}[t]{\wd\naastbox}\vspace{0pt}\usebox{\naastbox}\end{minipage}\par\addvspace{1mm}}
% Een rij figuren, gecentreerd
\newcommand{\figuren}[1]{\par\addvspace{1mm}{\centering #1\par}\addvspace{1.5mm}}
% Opgave met een figuur rechts: \begin{metgrafiek}{tikz} opgave en deelvragen \end{metgrafiek}
\newsavebox{\grafiekbox}
\newenvironment{metgrafiek}[2][]{%
  \sbox{\grafiekbox}{#2}%
  \begin{minipage}[t]{\dimexpr\linewidth-\wd\grafiekbox-3mm\relax}%
}{%
  \end{minipage}\hfill
  \begin{minipage}[t]{\wd\grafiekbox}\vspace{-2mm}\usebox{\grafiekbox}\end{minipage}\par}

%----------------------------------------------------------------------------------------
% TABELLEN
%----------------------------------------------------------------------------------------
% Tabel met waarden, horizontaal en met een raster: \waardetabel{aantal kolommen}{x-waarden}{y-waarden},
% de waarden gescheiden door &, in wiskundemodus. Met eigen namen: \waardetabel[$t$ (min)][$T$ (°C)]{5}{...}{...}
\NewDocumentCommand{\waardetabel}{O{$x$} O{$y$} m m m}{%
  {\renewcommand{\arraystretch}{1.35}%
  \begin{tabular}{|>{\columncolor{accent!12}}c|*{#3}{>{$}c<{$}|}}\hline #1 & #4 \\ \hline #2 & #5 \\ \hline\end{tabular}}}

%----------------------------------------------------------------------------------------
% GRAFIEKEN, binnen een tikzpicture. Zet de schaal met x=... en y=... (bv. [x=.5cm,y=.5cm]).
% Zoals in Functies.
%----------------------------------------------------------------------------------------
% Rooster met assen: \rooster{xmin}{xstap}{xmax}{ymin}{ystap}{ymax}, met getallen bij elke lijn (niet bij 0).
% Andere namen voor de assen: \rooster[$t$][$h$]{...}. \roosterleeg doet hetzelfde zonder getallen.
% Getallen bij een grotere stap: \roosterleeg, en dan \xgetallen{2,4,6} en \ygetallen{10,20}.
\NewDocumentCommand{\roosterleeg}{O{$x$} O{$y$} m m m m m m}{%
  \draw[black!14,very thin,xstep=#4,ystep=#7] (#3,#6) grid (#5,#8);
  \draw[->,grijs] (#3,0) -- ({#5+.45*#4},0) node[right,text=black,font=\footnotesize] {#1};
  \draw[->,grijs] (0,#6) -- (0,{#8+.45*#7}) node[above,text=black,font=\footnotesize] {#2};}
\NewDocumentCommand{\rooster}{O{$x$} O{$y$} m m m m m m}{%
  \roosterleeg[#1][#2]{#3}{#4}{#5}{#6}{#7}{#8}
  \pgfmathsetmacro{\rtweex}{#3+#4}\pgfmathsetmacro{\rtweey}{#6+#7}
  \foreach \t in {#3,\rtweex,...,#5}{\pgfmathtruncatemacro{\tt}{round(\t)}\ifdim\t pt=0pt\else\node[below,inner sep=1.5pt,text=black,font=\scriptsize] at (\t,0) {$\tt$};\fi}
  \foreach \t in {#6,\rtweey,...,#8}{\pgfmathtruncatemacro{\tt}{round(\t)}\ifdim\t pt=0pt\else\node[left,inner sep=1.5pt,text=black,font=\scriptsize] at (0,\t) {$\tt$};\fi}}
\newcommand{\xgetallen}[1]{\foreach \t in {#1}{\node[below,inner sep=1.5pt,text=black,font=\scriptsize] at (\t,0) {$\t$};}}
\newcommand{\ygetallen}[1]{\foreach \t in {#1}{\node[left,inner sep=1.5pt,text=black,font=\scriptsize] at (0,\t) {$\t$};}}
\newcommand{\xlabels}[1]{\foreach \t/\l in {#1}{\node[below,inner sep=1.5pt,text=black,font=\scriptsize] at (\t,0) {$\l$};}}
\newcommand{\ylabels}[1]{\foreach \t/\l in {#1}{\node[left,inner sep=1.5pt,text=black,font=\scriptsize] at (0,\t) {$\l$};}}
% De oorsprong: \oorsprong zet een 0 linksonder de oorsprong
\newcommand{\oorsprong}{\node[below left,inner sep=1.5pt,text=black,font=\scriptsize] at (0,0) {$0$};}

% Een rechte y = mx + q: \rechte[opties]{xmin}{xmax}{m}{q}
\newcommand{\rechte}[5][]{\draw[accent,thick,#1] (#2,{(#4)*(#2)+(#5)}) -- (#3,{(#4)*(#3)+(#5)});}
% Een rechte ax + by = c (met b verschillend van 0): \rechteabc[opties]{xmin}{xmax}{a}{b}{c}
\newcommand{\rechteabc}[6][]{\draw[accent,thick,#1] (#2,{((#6)-(#4)*(#2))/(#5)}) -- (#3,{((#6)-(#4)*(#3))/(#5)});}
% Een verticale rechte x = k: \verticaal[opties]{k}{ymin}{ymax}
\newcommand{\verticaal}[4][]{\draw[accent,thick,#1] (#2,#3) -- (#2,#4);}
% Een grafiek: \kromme[opties]{xmin}{xmax}{voorschrift in \x}
\newcommand{\kromme}[4][]{\draw[accent,thick,#1] plot[domain=#2:#3,samples=90,smooth] (\x,{#4});}
% Een punt: \punt{x}{y}, een open bolletje: \openpunt{x}{y}; kleur met de optie
\newcommand{\punt}[3][accent]{\fill[#1] ({#2},{#3}) circle (2.2pt);}
\newcommand{\openpunt}[3][accent]{\filldraw[#1,fill=white,thick] ({#2},{#3}) circle (2.2pt);}
% Stippellijnen om een punt af te lezen: \aflezen{x}{y} (van de x-as naar het punt en naar de y-as)
\newcommand{\aflezen}[3][grijs]{\draw[#1,dashed,thin] ({#2},0) -- ({#2},{#3}) -- (0,{#3});}
% Het trapje van de richtingscoëfficiënt: \trapje{x0}{y0}{dx}{dy}{label dx}{label dy}
\newcommand{\trapje}[6]{%
  \draw[inkt,thick,dashed,->,>=stealth] (#1,#2) -- ({#1+#3},#2);
  \draw[inkt,thick,dashed,->,>=stealth] ({#1+#3},#2) -- ({#1+#3},{#2+#4});
  \node[font=\scriptsize,text=inkt,anchor=\ifdim#4pt<0pt south\else north\fi,inner sep=2pt] at ({#1+#3/2},#2) {#5};
  \node[font=\scriptsize,text=inkt,anchor=\ifdim#3pt<0pt east\else west\fi,inner sep=2pt] at ({#1+#3},{#2+#4/2}) {#6};}
% Een tweede rechte in dezelfde figuur: stijl tweede (gestreept, donker), en hulplijn
\tikzset{tweede/.style={inkt,thick,dashed},hulplijn/.style={grijs,thick,densely dotted}}

%----------------------------------------------------------------------------------------
% De balans met twee soorten onbekenden (binnen een tikzpicture), zoals in Eerstegraadsvergelijkingen.
% \balansxy{zakjes links}{bollen links}{blokjes links}{zakjes rechts}{bollen rechts}{blokjes rechts}
% Een zakje (getint vierkant, met x) weegt x, een bol (wit, met y) weegt y, een blokje (grijs) weegt 1.
% De blokjes staan in kolommen van drie hoog. Een schaal is 3,9 breed: samen hoogstens 4 zakjes en bollen
% met 3 blokjes, of 3 met 6 blokjes. De hele balans is 9,8 eenheden breed.
%----------------------------------------------------------------------------------------
\newcommand{\balansschaalxy}[4]{%
  \draw[grijs,line width=1pt] (#1,0) -- (#1,.6);
  \draw[grijs,line width=1.4pt] ({#1-1.95},.6) -- ({#1+1.95},.6);
  \pgfmathtruncatemacro{\bkol}{ceil((#4)/3)}%
  \pgfmathsetmacro{\bstart}{#1-(((#2)+(#3))*.78+\bkol*.46-.08)/2}%
  \ifnum#2>0
    \foreach \i in {1,...,#2}{%
      \pgfmathsetmacro{\bx}{\bstart+(\i-1)*.78}%
      \draw[accent,fill=accent!22,rounded corners=1pt] (\bx,.67) rectangle ++(.7,.7);
      \node[font=\scriptsize,text=accent,inner sep=0pt] at ({\bx+.35},1.02) {$x$};}%
  \fi
  \ifnum#3>0
    \foreach \i in {1,...,#3}{%
      \pgfmathsetmacro{\bx}{\bstart+(#2)*.78+(\i-1)*.78}%
      \draw[inkt,fill=white,thick] ({\bx+.35},1.02) circle (.35);
      \node[font=\scriptsize,text=inkt,inner sep=0pt] at ({\bx+.35},1.02) {$y$};}%
  \fi
  \ifnum#4>0
    \foreach \i in {1,...,#4}{%
      \pgfmathtruncatemacro{\bk}{floor((\i-1)/3)}%
      \pgfmathtruncatemacro{\br}{mod(\i-1,3)}%
      \pgfmathsetmacro{\bx}{\bstart+((#2)+(#3))*.78+\bk*.46}%
      \pgfmathsetmacro{\by}{.67+\br*.46}%
      \draw[grijs,fill=black!12] (\bx,\by) rectangle ++(.38,.38);}%
  \fi}
\newcommand{\balansxy}[6]{%
  \fill[grijs] (-.7,-1.3) -- (.7,-1.3) -- (0,-.8) -- cycle;
  \draw[grijs,line width=1.3pt] (0,-1) -- (0,0);
  \draw[grijs,line width=1.8pt] (-2.95,0) -- (2.95,0);
  \fill[grijs] (0,0) circle (2.2pt);
  \balansschaalxy{-2.95}{#1}{#2}{#3}
  \balansschaalxy{2.95}{#4}{#5}{#6}}
